%% kapa: exam=1 round=30 year=2019 period=2 form=A subject=accounting subject_ko=회계학 seq=13 number=53
%% source: Q-Net 감정평가사 기출문제 (cst00309) artlSeq=5205205 file=extracted/1008/5205205/2019년 제30회 감정평가사 1차시험 문제지/2019년 제30회 감정평가사 1차시험 문제지/2019년도 제30회 감정평가사 1차 2교시 A형.hwp
%% conversion: kordoc 4.24.51 -> tools/gen1.py (auto); review=proofread-sonnet; flags=none
\begin{question}{53}
(주)감평은 리스이용자로 사무실용 건물을 20x1년 초부터 4년간 리스하는 계약(연간리스료 매년 말 ￦90,000 지급)을 체결하였다. (주)감평은 리스개시일인 20x1년 초에 리스부채로 ￦311,859을 인식하였다. 한편, 2년이 경과된 20x3년 초 (주)감평은 리스회사와 매년 말 연간 리스료 ￦70,000을 지급하기로 합의하였다. 20x3년 초 리스변경을 반영한 후 (주)감평의 리스부채 장부금액은? (단, 리스의 내재이자율은 쉽게 산정할 수 없으나, 리스개시일과 20x3년 초 리스이용자인 (주)감평의 증분차입이자율은 각각 연 6 \%와 연 8 \%이다.) \par
{\small\centering
\begin{tabular}{|c|c|c|}
\hline
\multirow{2}{*}{\makecell{기간}} & \multicolumn{2}{|c|}{\makecell{정상연금 ￦1의 현재가치}} \\
\cline{2-3}
 & \makecell{6 \%} & \makecell{8 \%} \\
\hline
\makecell{1} & \makecell{0.9434} & \makecell{0.9259} \\
\hline
\makecell{2} & \makecell{1.8334} & \makecell{1.7833} \\
\hline
\makecell{3} & \makecell{2.6730} & \makecell{2.5771} \\
\hline
\makecell{4} & \makecell{3.4651} & \makecell{3.3121} \\
\hline
\end{tabular}\par}

\begin{choices}
  \item ￦124,831
  \item ￦128,338
  \item ￦159,456
  \item ￦231,847
  \item ￦242,557
\end{choices}
\end{question}
\answer{1}
